% SEGMENT OLP-0644-S01
% Part: first-order-logic
% Chapter: completeness
% Section: maximally-consistent-sets

% அதிகபட்ச முரணற்ற கணங்களின் வரையறை. நிறைவு நிறுவலுக்குத் தேவையான
% பண்புகள், பயன்படுத்தப்படும் நிறுவல் முறையின் அத்தியாயத்தில் உள்ளன.

\documentclass[../../../include/open-logic-section]{subfiles}

\begin{document}

\olfileid{fol}{com}{mcs}
\olsection{அதிகபட்ச முரணற்ற \usetoken{P}{sentence}-கணங்கள்}

\begin{defn}[அதிகபட்ச முரணற்ற கணம்]
\ollabel{mcs}
!!{sentence}s கொண்ட கணம்~$\Gamma$, பின்வரும் இரண்டு நிபந்தனைகளும்
நிறைவேறினால், அப்போதும் அப்போது மட்டுமே \emph{அதிகபட்ச முரணற்றது}:
\begin{enumerate}
\item $\Gamma$ முரணற்றது; மேலும்
\item $\Gamma \subsetneq \Gamma'$ எனில், $\Gamma'$ முரணுடையது.
\end{enumerate}
\end{defn}

% SEGMENT OLP-0644-S02
\begin{explain}
மேலுள்ளதற்குச் சமானமான மற்றொரு வரையறை இதுதான்: வாக்கியக்
கணம்~$\Gamma$, பின்வரும் இரண்டு நிபந்தனைகளும் நிறைவேறினால்,
அப்போதும் அப்போது மட்டுமே \emph{அதிகபட்ச முரணற்றது}:
\begin{enumerate}
\item $\Gamma$ முரணற்றது; மேலும்
\item $\Gamma \cup \{ !A \}$ முரணற்றது எனில், $!A \in \Gamma$.
\end{enumerate}
வேறு சொற்களில், ஏற்கெனவே~$\Gamma$ இல் இல்லாத !!{sentence}s ஐ
அதிகபட்ச முரணற்ற கணம்~$\Gamma$ உடன் சேர்த்தால், அதன் விளைவான பெரிய கணம்
முரணுடையதாகிவிடும்.
\end{explain}

% SEGMENT OLP-0644-S03
\begin{explain}
நிறைவு நிறுவலில் அதிகபட்ச முரணற்ற கணங்கள் முக்கியமானவை. ஏனெனில்,
!!{sentence}s கொண்ட எந்த முரணற்ற கணம்~$\Gamma$ உம் ஓர் அதிகபட்ச
முரணற்ற கணம்~$\Gamma^*$ இல் அடங்கும் என்பதை உறுதிசெய்யலாம். மேலும்,
அதிகபட்ச முரணற்ற கணம், ஒவ்வொரு !!{sentence}~$!A$ க்கும் $!A$
அல்லது அதன் மறுப்பு~$\lnot !A$ இவற்றில் ஒன்றைக் கொண்டிருக்கும்.
இது குறிப்பாக அணு !!{sentence}s க்கும் பொருந்துகிறது. எனவே
!!{constant}s ஐ ஏற்றவாறு சேர்த்து விரிவாக்கிய மொழியில் உள்ள
அதிகபட்ச முரணற்ற கணத்திலிருந்து, எந்த அணு !!{sentence}s ஆனவை
$\Gamma^*$ இல் உள்ளன என்பதன்படி !!{predicate}s இன் பொருள்விளக்கம்
வரையறுக்கப்படும் !!a{structure} ஒன்றைக் கட்டமைக்கலாம். இந்த
!!{structure} ஆனது~$\Gamma^*$ இல் உள்ள அனைத்து !!{sentence}s ஐயும்
(ஆகவே~$\Gamma$ இல் உள்ள அனைத்தையும்) மெய்யாக்கும் என்று பின்னர்
காட்டலாம். இந்தப் பின்னைய கூற்றை நிறுவ, $\lnot !A \in
\Gamma^*$ அப்போதும் அப்போது மட்டுமே $!A \notin \Gamma^*$;
$(!A \lor !B) \in \Gamma^*$ அப்போதும் அப்போது மட்டுமே $!A
\in \Gamma^*$ அல்லது $!B \in \Gamma^*$; இதுபோன்ற பண்புகள் தேவை.
\end{explain}

% SEGMENT OLP-0644-S04
\begin{prop}
\ollabel{prop:mcs}
$\Gamma$ அதிகபட்ச முரணற்றது என்க. அப்போது:
\begin{enumerate}
\item \ollabel{prop:mcs-prov-in} $\Gamma \Proves !A$ எனில், $!A \in
  \Gamma$.

\item \ollabel{prop:mcs-either-or} எந்த $!A$ க்கும், $!A \in
  \Gamma$ அல்லது $\lnot !A \in \Gamma$.

\tagitem{prvAnd}{\ollabel{prop:mcs-and} $(!A \land !B) \in \Gamma$
  அப்போதும் அப்போது மட்டுமே $!A \in \Gamma$, $!B \in \Gamma$ இரண்டும்.}{}

\tagitem{prvOr}{\ollabel{prop:mcs-or} $(!A \lor !B) \in \Gamma$
  அப்போதும் அப்போது மட்டுமே $!A \in \Gamma$ அல்லது $!B \in \Gamma$.}{}

\tagitem{prvIf}{\ollabel{prop:mcs-if} $(!A \lif !B) \in \Gamma$
  அப்போதும் அப்போது மட்டுமே $!A \notin \Gamma$ அல்லது $!B \in \Gamma$.}{}
\end{enumerate}
\end{prop}

% SEGMENT OLP-0644-S05
\begin{proof}
பின்வரும் அனைத்திலும் $\Gamma$ அதிகபட்ச முரணற்றது எனக் கொள்வோம்.
\begin{enumerate}
\item $\Gamma \Proves !A$ எனில், $!A \in \Gamma$.

$\Gamma \Proves !A$ என்க. இதற்கு மாறாக $!A
\notin \Gamma$ எனக் கொள்வோம். $\Gamma$ அதிகபட்ச முரணற்றது என்பதால்,
$\Gamma \cup \{!A\}$ முரணுடையது; எனவே $\Gamma \cup \{!A\} \Proves
\lfalse$. பின்வரும் மேற்கோளின்படி,
\tagrefs{prfSC/{fol:seq:prv:prop:provability-contr},prfND/{fol:ntd:prv:prop:provability-contr}},
$\Gamma$ முரணுடையது. இது $\Gamma$ முரணற்றது என்ற கருதுகோளுக்கு
முரணாகும். ஆகவே $!A \notin
\Gamma$ என்பது இயலாது; எனவே $!A \in \Gamma$.

\item எந்த $!A$ க்கும், $!A \in \Gamma$ அல்லது $\lnot !A \in \Gamma$.

இதற்கு மாறாக, ஏதோ~$!A$ க்கு $!A \notin \Gamma$ மற்றும்
$\lnot !A \notin \Gamma$ இரண்டும் எனக் கொள்வோம். $\Gamma$ அதிகபட்ச
முரணற்றது என்பதால், $\Gamma \cup \{!A\}$ மற்றும்
$\Gamma \cup \{\lnot !A\}$ இரண்டும் முரணுடையவை. எனவே
$\Gamma \cup \{!A\} \Proves \lfalse$ மற்றும் $\Gamma \cup
\{\lnot !A\} \Proves \lfalse$. பின்வரும் மேற்கோளின்படி,
\tagrefs{prfSC/{fol:seq:prv:prop:provability-exhaustive},prfND/{fol:ntd:prv:prop:provability-exhaustive}},
$\Gamma$ முரணுடையது; இது முரண்பாடு. ஆகவே அத்தகைய
!!a{sentence}~$!A$ இருக்க முடியாது. இதனால் ஒவ்வொரு $!A$ க்கும்,
$!A \in \Gamma$ அல்லது $\lnot !A
\in \Gamma$.

% SEGMENT OLP-0644-S06
\tagitem{defAnd}{}{%
\iftag{probAnd}{பயிற்சி.}{%
$(!A \land !B) \in \Gamma$ அப்போதும் அப்போது மட்டுமே
$!A \in \Gamma$, $!B \in \Gamma$ இரண்டும்:

முன்னோக்கிய திசைக்கு, $(!A \land !B) \in \Gamma$ என்க. அப்போது
$\Gamma \Proves !A \land !B$. பின்வரும் மேற்கோளின்படி,
\tagrefs{prfSC/{fol:seq:ppr:prop:provability-land-left},prfND/{fol:ntd:ppr:prop:provability-land-left}},
$\Gamma \Proves !A$ மற்றும் $\Gamma \Proves !B$. தேவையானவாறு,
\olref{prop:mcs-prov-in} படி $!A \in \Gamma$ மற்றும் $!B \in \Gamma$.

மறுதிசைக்கு, $!A \in \Gamma$ மற்றும் $!B \in
\Gamma$ என்க. அப்போது $\Gamma \Proves !A$ மற்றும்
$\Gamma \Proves !B$. பின்வரும் மேற்கோளின்படி,
\tagrefs{prfSC/{fol:seq:ppr:prop:provability-land-right},prfND/{fol:ntd:ppr:prop:provability-land-right}},
$\Gamma \Proves !A \land !B$. \olref{prop:mcs-prov-in} படி,
$(!A \land !B) \in \Gamma$.}}

% SEGMENT OLP-0644-S07
\tagitem{defOr}{}{%
\iftag{probOr}{பயிற்சி.}{%
$(!A \lor !B) \in \Gamma$ அப்போதும் அப்போது மட்டுமே
$!A \in \Gamma$ அல்லது $!B \in \Gamma$.

முன்னோக்கிய திசையின் எதிர்மறைநேர்மாற்றை நிறுவ, $!A
\notin \Gamma$ மற்றும் $!B \notin \Gamma$ என்க. நாம் $(!A \lor
!B) \notin \Gamma$ என்பதைக் காட்ட வேண்டும். $\Gamma$ அதிகபட்ச
முரணற்றது என்பதால், $\Gamma \cup \{!A\} \Proves \lfalse$ மற்றும்
$\Gamma \cup \{!B\} \Proves \lfalse$. பின்வரும் மேற்கோளின்படி,
\tagrefs{prfSC/{fol:seq:ppr:prop:provability-lor},prfND/{fol:ntd:ppr:prop:provability-lor}},
$\Gamma \cup \{(!A \lor !B)\}$ முரணுடையது. ஆகவே தேவையானபடி,
$(!A \lor !B)
\notin \Gamma$.

மறுதிசைக்கு, $!A \in \Gamma$ அல்லது $!B \in
\Gamma$ என்க. அப்போது $\Gamma \Proves !A$ அல்லது
$\Gamma \Proves !B$. பின்வரும் மேற்கோளின்படி,
\tagrefs{prfSC/{fol:seq:ppr:prop:provability-lor},prfND/{fol:ntd:ppr:prop:provability-lor}},
$\Gamma \Proves !A \lor !B$. \olref{prop:mcs-prov-in} படி,
தேவையானபடி $(!A \lor
!B) \in \Gamma$.}}

% SEGMENT OLP-0644-S08
\tagitem{defIf}{}{%
\iftag{probIf}{பயிற்சி.}{%
$(!A \lif !B) \in \Gamma$ அப்போதும் அப்போது மட்டுமே
$!A \notin \Gamma$ அல்லது $!B \in \Gamma$:

முன்னோக்கிய திசைக்கு, $(!A \lif !B) \in \Gamma$ என்க. இதற்கு
மாறாக $!A \in \Gamma$ மற்றும் $!B \notin \Gamma$ எனக் கொள்வோம்.
இந்தக் கருதுகோள்களின்படி $\Gamma \Proves !A \lif !B$ மற்றும்
$\Gamma \Proves !A$. பின்வரும் மேற்கோளின்படி,
\tagrefs{prfSC/{fol:seq:ppr:prop:provability-lif-left},prfND/{fol:ntd:ppr:prop:provability-lif-left}},
$\Gamma \Proves !B$. ஆனால் \olref{prop:mcs-prov-in} படி $!B \in
\Gamma$; இது $!B \notin \Gamma$ என்ற கருதுகோளுக்கு முரணாகும்.

மறுதிசையில், முதலில் $!A \notin
\Gamma$ என்ற நிலையை எடுத்துக்கொள்வோம். \olref{prop:mcs-either-or}
படி $\lnot !A \in \Gamma$; ஆகவே $\Gamma \Proves \lnot !A$.
பின்வரும் மேற்கோளின்படி,
\tagrefs{prfSC/{fol:seq:ppr:prop:provability-lif-right},prfND/{fol:ntd:ppr:prop:provability-lif-right}},
$\Gamma \Proves !A \lif !B$. மீண்டும் \olref{prop:mcs-prov-in}
படி, தேவையானபடி $(!A \lif !B) \in \Gamma$.

இப்போது $!B \in \Gamma$ என்ற நிலையை எடுத்துக்கொள்வோம். அப்போது
$\Gamma \Proves !B$. பின்வரும் மேற்கோளின்படி,
\tagrefs{prfSC/{fol:seq:ppr:prop:provability-lif-right},prfND/{fol:ntd:ppr:prop:provability-lif-right}},
$\Gamma \Proves !A \lif !B$. \olref{prop:mcs-prov-in} படி,
$(!A \lif
!B) \in \Gamma$.}}
\end{enumerate}
\end{proof}

% SEGMENT OLP-0644-S09
\begin{prob}
\olref[fol][com][mcs]{prop:mcs} இன் நிறுவலை நிறைவு செய்க.
\end{prob}

\end{document}
